% Provenance (never printed): NSI terminale "Algorithmique: recherche textuelle (algorithme de Boyer-Moore, ici
% sous la forme simplifiee de Horspool; interet du pretraitement du motif; cout non exigible)". See
% sources/book2/BRIEFS.md.
\chapter{Searching in Text}\label{ch:g12:text-search}

A text editor finds a word in a book of a million letters before the finger leaves the key. The obvious way is to
try the word at every position of the text, comparing letters until one differs. It works, and it reads every
letter of the text at least once. A cleverer search looks at the word before it starts, and then, at most positions,
reads one letter and jumps several places at once.

This chapter writes both searches, counts the letters they compare, and shows when the jumps are long and when they
are short. The cost is measured on generated texts rather than proved: the guarantee of the cleverer method is that
it never misses an \omterm{def:g12:text-search:pattern}{occurrence}, and that much is proved.

\begin{recall}
A string is a sequence of characters, indexed from $0$; \code{s[i:j]} is the slice from $i$ up to $j$ (Book 1, grade
7; chapter \ref{ch:g10:text-encodings}). The cost of an algorithm is counted in elementary steps, as a function of the
size of its input (chapter \ref{ch:g11:cost-of-an-algorithm}). A dictionary maps keys to values, and \code{d.get(k,
default)} gives a value with a default for absent keys (Book 1, grade 8).
\end{recall}

\section{The naive search}\label{sec:g12:text-search:naive}

\begin{definition}[Pattern, occurrence]\label{def:g12:text-search:pattern}
In text search, the \emph{pattern}\index{pattern (text search)} is the string looked for, of length $m$, in a
\emph{text} of length $n$. An \emph{occurrence}\index{occurrence} of the pattern is a position $i$ of the text with
\code{text[i:i + m] == p}. Placing the pattern under the text from position $i$ is called an \emph{alignment}; there
are $n - m + 1$ of them.
\end{definition}

The naive search tries every alignment, from left to right. At each one it compares the letters of the \omterm{def:g12:text-search:pattern}{pattern} with
those of the text above them, from the first, and gives up at the first difference.

\omcode{code/grade-12/08-text-search/listings/naive.py}{1}{24}{The naive search, with its comparisons
counted.}\label{lst:g12:text-search:naive}

\omoutput{code/grade-12/08-text-search/outputs/naive.txt}{Positions found, and comparisons made, in a text of $27$
letters.}\label{out:g12:text-search:naive}

\begin{omfigure}
\begin{tikzpicture}[font=\ttfamily\small, x=4.4mm, y=-5.2mm]
  \foreach \c [count=\k from 0] in {a,{ },m,a,n,{ },a,{ },p,l,a,n} \node at (\k,0) {\c};
  \foreach \k in {0,...,11} \node[font=\tiny, black!50] at (\k,-0.8) {\k};
  \foreach \i/\r in {0/1, 1/2, 2/3, 3/4} {
    \foreach \c [count=\j from 0] in {p,l,a,n} {
      \pgfmathtruncatemacro{\k}{\i+\j}
      \ifnum\j=0 \node[draw, fill=omThm!15, minimum size=4mm, inner sep=0pt] at (\k,\r) {\c};
      \else \node[black!35] at (\k,\r) {\c}; \fi
    }
  }
  \foreach \c [count=\j from 0] in {p,l,a,n} \node[draw, fill=omMeth!20, minimum size=4mm, inner sep=0pt] at (8+\j,5.4) {\c};
  \node[font=\small\rmfamily, black!60] at (8,3.2) {$\vdots$};
  \node[font=\footnotesize\rmfamily, right, align=left] at (12.6,2.5) {each alignment fails\\ at its first letter: one\\ comparison};
  \node[font=\footnotesize\rmfamily, right] at (12.6,5.4) {position $8$: four comparisons, found};
\end{tikzpicture}

{\small The naive search for \code{plan}. Each row is one alignment; the boxed letter is the one compared, the grey
letters are never looked at. In ordinary text most alignments fail at the first or second letter, so the search makes
a little more than one comparison per position.}
\end{omfigure}

In the worst case, every alignment matches all letters but the last. In a text of $n$ letters \code{a}, the \omterm{def:g12:text-search:pattern}{pattern}
\code{aa...ab} costs $m$ comparisons at each of the $n - m + 1$ alignments.

\omcode{code/grade-12/08-text-search/listings/worst_naive.py}{1}{9}{The worst case of the naive
search.}\label{lst:g12:text-search:worst}

\omoutput{code/grade-12/08-text-search/outputs/worst_naive.txt}{Exactly $(n - m + 1) \times m$
comparisons.}\label{out:g12:text-search:worst}

\section{Looking at the pattern first}\label{sec:g12:text-search:preprocessing}

Suppose the \omterm{def:g12:text-search:pattern}{pattern} is \code{STATION} and, at some alignment, the text letter under its last place is a \code{K}. No
letter of \code{STATION} is a \code{K}. Any alignment that puts a letter of the \omterm{def:g12:text-search:pattern}{pattern} under that \code{K} fails; the
\omterm{def:g12:text-search:pattern}{pattern} can jump past it, seven places at once. If the letter were a \code{T}, the \omterm{def:g12:text-search:pattern}{pattern} could jump until its own
last \code{T}, not counting the final place, comes under it: three places.

These distances depend on the \omterm{def:g12:text-search:pattern}{pattern} only, so they can be computed once, before reading the text. Doing work on
the \omterm{def:g12:text-search:pattern}{pattern} in advance to save work on the text is called \emph{preprocessing}.

\begin{definition}[Shift table]\label{def:g12:text-search:shift-table}
The \emph{shift table}\index{shift table} of a \omterm{def:g12:text-search:pattern}{pattern} $p$ of length $m$ gives, for each letter $c$, the distance from
the last place of $c$ in $p$, the final place excluded, to the final place: $m - 1 - j$, where $j$ is that last place.
A letter absent from the first $m - 1$ places of $p$ gets the shift $m$.
\end{definition}

\omcode{code/grade-12/08-text-search/listings/shift_table.py}{1}{15}{The \omterm{def:g12:text-search:shift-table}{shift table}, as a dictionary; absent letters
are left out.}\label{lst:g12:text-search:shift}

\omoutput{code/grade-12/08-text-search/outputs/shift_table.txt}{The later \code{T} of \code{STATION} overwrote the
earlier one.}\label{out:g12:text-search:shift}

\begin{omfigure}
\begin{tikzpicture}[font=\ttfamily\small, x=7mm, y=-6mm]
  \foreach \c [count=\j from 0] in {S,T,A,T,I,O,N} {
    \node[draw, minimum width=7mm, minimum height=6mm, fill=omDef!8] at (\j,0) {\c};
    \node[font=\tiny\rmfamily, black!50] at (\j,-0.85) {\j};
  }
  \foreach \c/\s [count=\k from 0] in {S/6, T/3, A/4, I/2, O/1, {other}/7} {
    \node[draw, minimum width=11mm, minimum height=6mm, fill=black!5] at (1.6*\k-0.4,2) {\c};
    \node[draw, minimum width=11mm, minimum height=6mm] at (1.6*\k-0.4,3) {\s};
  }
  \node[draw, minimum width=7mm, minimum height=6mm, fill=black!8] at (6,0) {N};
  \node[left, font=\footnotesize\rmfamily] at (-1.1,2) {letter}; \node[left, font=\footnotesize\rmfamily] at (-1.1,3) {shift};
  \draw[<->, omThm, thick] (3,0.55) -- node[below, font=\scriptsize\rmfamily] {$3$} (6,0.55);
\end{tikzpicture}

{\small The \omterm{def:g12:text-search:shift-table}{shift table} of \code{STATION} ($m = 7$). The last \code{T} before the final place is at $3$, three places
from the end: shift $3$. \code{N}, which appears only in the final place (shaded), and every letter absent from the
\omterm{def:g12:text-search:pattern}{pattern} get $7$.}
\end{omfigure}

\section{The Boyer--Moore--Horspool search}\label{sec:g12:text-search:horspool}

\begin{method}[Boyer--Moore--Horspool]\label{met:g12:text-search:horspool}
The \emph{Boyer--Moore--Horspool algorithm}\index{Boyer--Moore--Horspool algorithm} computes the \omterm{def:g12:text-search:shift-table}{shift table} of the
\omterm{def:g12:text-search:pattern}{pattern}, then places the \omterm{def:g12:text-search:pattern}{pattern} at position $i = 0$ and repeats while $i \le n - m$: compare the \omterm{def:g12:text-search:pattern}{pattern} with the text
above it from right to left, and record $i$ if all $m$ letters agree; then, whatever happened, read the text letter
$c$ under the \omterm{def:g12:text-search:pattern}{pattern}'s last place and move the \omterm{def:g12:text-search:pattern}{pattern} $\mathrm{shift}(c)$ places to the right.
\end{method}

\omcode{code/grade-12/08-text-search/listings/horspool.py}{1}{30}{Horspool's search, with a trace of the
alignments.}\label{lst:g12:text-search:horspool}

\omoutput{code/grade-12/08-text-search/outputs/horspool.txt}{Seven alignments and $13$ comparisons to find
\code{station} in $45$ letters; then \code{the}, three times.}\label{out:g12:text-search:horspool}

\begin{omfigure}
\begin{tikzpicture}[font=\ttfamily\small, x=4.6mm, y=-5.4mm]
  \foreach \c [count=\k from 0] in {t,h,e,{ },b,u,s,{ },t,o,{ },t,h} \node at (\k,0) {\c};
  \foreach \k in {0,...,12} \node[font=\tiny, black!50] at (\k,-0.8) {\k};
  \foreach \c [count=\j from 0] in {s,t,a,t,i,o} \node[black!45] at (\j,1.3) {\c};
  \node[draw, fill=omThm!15, minimum size=4mm, inner sep=0pt] at (6,1.3) {n};
  \draw[omThm, thick] (5.6,-0.35) rectangle (6.4,0.35);
  \node[font=\footnotesize\rmfamily, right, align=left] at (13.2,0.6) {\code{s} $\ne$ \code{n}: no;\\ under the last place: \code{s}};
  \foreach \c [count=\j from 0] in {s,t,a,t,i,o,n} \node[draw, fill=omMeth!12, minimum size=4mm, inner sep=0pt] at (6+\j,3.2) {\c};
  \draw[->, thick, omThm] (0,1.9) to[bend right=12] node[below left, font=\scriptsize\rmfamily] {jump $6$} (6,2.7);
  \draw[densely dotted, omThm] (6,0.4) -- (6,2.8);
\end{tikzpicture}

{\small A jump of Horspool's search for \code{station}. At position $0$, the first step of the trace, the text letter under the last
place is \code{s}; the \omterm{def:g12:text-search:pattern}{pattern}'s only \code{s} is $6$ places from its end, so the \omterm{def:g12:text-search:pattern}{pattern} jumps $6$ places, bringing its
\code{s} under the text's \code{s}. Any smaller jump would put some other letter of the \omterm{def:g12:text-search:pattern}{pattern} under that \code{s}.}
\end{omfigure}

Why are no \omterm{def:g12:text-search:pattern}{occurrences} skipped? Let $c$ be the text letter under the last place, at alignment $i$, and $s$ its shift.
For a jump of $d$ places with $1 \le d < s$, the letter of the \omterm{def:g12:text-search:pattern}{pattern} that would come under $c$ is at place
$m - 1 - d$, which is not the last place and lies after the last place of $c$ among the first $m - 1$; so it is not
$c$, and that alignment fails. The alignments jumped over all fail, and the search tries every one that could succeed
(exercise \ref{exo:g12:text-search:12}).

\begin{history}[Boyer and Moore, 1977; Horspool, 1980]
Robert S.\ Boyer and J Strother Moore published ``A fast string searching algorithm'' in the \emph{Communications of
the ACM} in 1977; their method uses the \omterm{def:g12:text-search:shift-table}{shift table} of this chapter and a second, more elaborate table. In 1980,
R.\ Nigel Horspool showed in ``Practical fast searching in strings'' that the first table alone gives most of the
speed.
% ledger: g12:boyer-moore-1977, g12:horspool-1980
\end{history}

\section{The cost in practice}\label{sec:g12:text-search:cost}

How many comparisons does each search make for each letter of the text? The answer depends on the text. The program
draws random texts of $100\,000$ letters, from $26$ letters and from four (\code{A}, \code{C}, \code{G}, \code{T}, as
in genetic sequences), and cuts the \omterm{def:g12:text-search:pattern}{patterns} from the texts so that each occurs at least once.

\omcode{code/grade-12/08-text-search/listings/compare_counts.py}{1}{19}{Comparisons per text letter on generated
texts.}\label{lst:g12:text-search:compare}

\omoutput{code/grade-12/08-text-search/outputs/compare_counts.txt}{With $26$ letters, Horspool's count falls as the
\omterm{def:g12:text-search:pattern}{pattern} grows; with four, it does not go on falling.}\label{out:g12:text-search:compare}

\begin{omfigure}
\begin{tikzpicture}
\begin{axis}[omrecord, width=11cm, height=6cm, xlabel={length $m$ of the pattern}, ylabel={comparisons per letter},
    ymin=0, ymax=2.1, legend pos=north east, legend style={font=\footnotesize}, legend columns=2]
  \addplot[omThm, thick, dashed] table[x=m, y=naive4]{code/grade-12/08-text-search/figures/per_letter.dat};
  \addplot[omThm, thick, mark=square*, mark size=1.2pt] table[x=m, y=hors4]{code/grade-12/08-text-search/figures/per_letter.dat};
  \addplot[omDef, thick, dashed] table[x=m, y=naive26]{code/grade-12/08-text-search/figures/per_letter.dat};
  \addplot[omDef, thick, mark=*, mark size=1.2pt] table[x=m, y=hors26]{code/grade-12/08-text-search/figures/per_letter.dat};
  \legend{naive{,} 4 letters, Horspool{,} 4 letters, naive{,} 26 letters, Horspool{,} 26 letters}
\end{axis}
\end{tikzpicture}

{\small Comparisons per text letter, averaged over five \omterm{def:g12:text-search:pattern}{patterns} for each length, on random texts of $50\,000$
letters, from \code{figures/per_letter.py}. The naive search stays near $1.04$ and $1.33$ whatever the \omterm{def:g12:text-search:pattern}{pattern}. With
$26$ letters Horspool's count falls roughly like $1/m$, to $0.06$ at $m = 32$; with four letters it stays around
$0.4$.}
\end{omfigure}

The jump is long only when the letter under the last place is rare in the \omterm{def:g12:text-search:pattern}{pattern}. With $26$ letters, most text
letters do not appear in a short \omterm{def:g12:text-search:pattern}{pattern} at all, and the \omterm{def:g12:text-search:pattern}{pattern} jumps its whole length. With four letters, every
letter appears among the last few places of any long \omterm{def:g12:text-search:pattern}{pattern}, so the jumps stay short, about four places, however
long the \omterm{def:g12:text-search:pattern}{pattern}. The worst case is as bad as the naive search's: in a text of \code{a}s, the \omterm{def:g12:text-search:pattern}{pattern}
\code{baa...a} costs $m$ comparisons at each alignment and jumps one place (exercise \ref{exo:g12:text-search:13}).

On the titles of the town library, joined into one text, the counts are those of ordinary text.

\omcode{code/grade-12/08-text-search/listings/town_titles.py}{1}{17}{Five words searched in the library's
titles.}\label{lst:g12:text-search:titles}

\omoutput{code/grade-12/08-text-search/outputs/town_titles.txt}{The longer the word, the fewer letters Horspool
reads; for \code{The}, three letters long, it still saves half.}\label{out:g12:text-search:titles}

\section{Searching in practice}\label{sec:g12:text-search:practice}

Counting comparisons is not the same as timing. A Python loop costs far more per step than the machine code of the
interpreter's own string methods.

\omcode{code/grade-12/08-text-search/listings/measure_find.py}{1}{31}{The two searches and \code{str.find}, timed on a
million letters.}\label{lst:g12:text-search:measure}

\omoutput{code/grade-12/08-text-search/outputs/measure_find.txt}{On the author's laptop. \code{str.find} answers in
under a millisecond, many times faster than either Python loop.}\label{out:g12:text-search:measure}

\begin{inpractice}[Searching in real programs]
Python's \code{str.find}, \code{in} and \code{str.count} are written in C and use a search of the Boyer--Moore family,
with a \omterm{def:g12:text-search:shift-table}{shift table} of their own. Text editors and command-line tools search the same way. Searching for many words at
once, such as every name of a list in a long text, uses a different preprocessing: the words are merged into one
structure that is followed letter by letter through the text, reading each letter once. It is studied in Book 4.
\end{inpractice}

\section{Exercises}\label{sec:g12:text-search:exercises}

\begin{exercise}[$\star$]\label{exo:g12:text-search:1}
Where does \code{ana} occur in \code{bananas}? Do the \omterm{def:g12:text-search:pattern}{occurrences} overlap?
\end{exercise}

\begin{exercise}[$\star$]\label{exo:g12:text-search:2}
Count the comparisons of \code{search_naive("abcabd", "abd")}.
\end{exercise}

\begin{exercise}[$\star$]\label{exo:g12:text-search:3}
Write the \omterm{def:g12:text-search:shift-table}{shift table} of \code{BANANA}.
\end{exercise}

\begin{exercise}[$\star$]\label{exo:g12:text-search:4}
List the alignments Horspool's search tries for \code{TOWN} in \code{THE OLD TOWN HALL}, and the jump after each.
\end{exercise}

\begin{exercise}[$\star$]\label{exo:g12:text-search:5}
Write \code{count_occurrences(text, p)} with slices, trying every position.
\end{exercise}

\begin{exercise}[$\star\star$]\label{exo:g12:text-search:6}
A student writes \code{for i in range(len(text) - len(p))}. Which \omterm{def:g12:text-search:pattern}{occurrences} are missed? Give the smallest test that
catches it.
\end{exercise}

\begin{exercise}[$\star\star$]\label{exo:g12:text-search:7}
Write \code{horspool_first(text, p)}, which returns the first \omterm{def:g12:text-search:pattern}{occurrence}, or $-1$, and stops as soon as it is found.
\end{exercise}

\begin{exercise}[$\star\star$]\label{exo:g12:text-search:8}
A student builds the \omterm{def:g12:text-search:shift-table}{shift table} with \code{range(m)} instead of \code{range(m - 1)}. What is the shift of \code{N}
for \code{TOWN}? What happens when the search meets an \code{N} under the last place?
\end{exercise}

\begin{exercise}[$\star\star$]\label{exo:g12:text-search:9}
For a text of $n$ letters, which \omterm{def:g12:text-search:pattern}{pattern} length $m$ makes the worst case of the naive search, $(n - m + 1)m$, the
largest? How many comparisons for $n = 1000$?
\end{exercise}

\begin{exercise}[$\star\star$]\label{exo:g12:text-search:10}
Where does \code{ABAB} occur in \code{ABABABAB}, by the naive search? What does \code{"ABABABAB".count("ABAB")} return,
and why the difference?
\end{exercise}

\begin{exercise}[$\star\star$]\label{exo:g12:text-search:11}
Write \code{search_ignore_case(text, p)}, which finds \code{TOWN} in \code{The Town and the town} twice.
\end{exercise}

\begin{exercise}[$\star\star\star$]\label{exo:g12:text-search:12}
Prove that Horspool's search finds every \omterm{def:g12:text-search:pattern}{occurrence}: no alignment it jumps over is an \omterm{def:g12:text-search:pattern}{occurrence}.
\end{exercise}

\begin{exercise}[$\star\star\star$]\label{exo:g12:text-search:13}
Show that for the text of $n$ letters \code{a} and the \omterm{def:g12:text-search:pattern}{pattern} \code{b} followed by $m - 1$ letters \code{a},
Horspool's search makes exactly $m(n - m + 1)$ comparisons.
\end{exercise}

\begin{exercise}[$\star\star\star$]\label{exo:g12:text-search:14}
Write \code{search_wildcard(text, p)}, in which a \code{?} of the \omterm{def:g12:text-search:pattern}{pattern} matches any one letter. Why does the \omterm{def:g12:text-search:shift-table}{shift
table} not work directly with a \code{?}?
\end{exercise}

\begin{exercise}[$\star\star\star$]\label{exo:g12:text-search:15}
On a random text of four letters, the mean jump of Horspool's search stays around four places, whatever the length of
the \omterm{def:g12:text-search:pattern}{pattern}. Explain why, and measure the mean jump for one \omterm{def:g12:text-search:pattern}{pattern} of length $32$.
\end{exercise}

\section{Problem: Searching a Long Chain of Letters}\label{sec:g12:text-search:problem}

\begin{problem}[{Weekend problem --- a million letters, four of them}]\label{pb:g12:text-search:1}
A genetic sequence is written with four letters, \code{A}, \code{C}, \code{G}, \code{T}. The program draws a random
chain of a million of them, cuts five \omterm{def:g12:text-search:pattern}{patterns} of each length from it, and runs both searches.

\omcode{code/grade-12/08-text-search/listings/chain.py}{1}{26}{Both searches on a million letters of four
kinds.}\label{lst:g12:text-search:chain}

\omoutput{code/grade-12/08-text-search/outputs/chain.txt}{Comparisons per letter, their ratio, and the mean jump of
the \omterm{def:g12:text-search:shift-table}{shift tables}.}\label{out:g12:text-search:chain}

\needspace{4\baselineskip}\medskip\noindent\textbf{Part I --- The chain.}
\begin{enumerate}
  \item Why are the \omterm{def:g12:text-search:pattern}{patterns} cut from the text rather than drawn at random?
  \item In a random chain, how often would you expect a given \omterm{def:g12:text-search:pattern}{pattern} of $4$ letters? Compare with the $378$ found in
    the $100\,000$ letters of \code{compare_counts.py}.
  \item Why is a \omterm{def:g12:text-search:pattern}{pattern} of $16$ letters very unlikely to occur twice?
  \item What does \code{assert found == found_h} check, and why is it in the program?
\end{enumerate}

\needspace{4\baselineskip}\medskip\noindent\textbf{Part II --- The naive search.}
\begin{enumerate}[resume]
  \item At an alignment, what is the chance that the first comparison fails? That the search reaches a second
    comparison? A third?
  \item Explain why the naive search makes about $1 + \frac14 + \frac1{16} + \dotsb$ comparisons per alignment, and
    why this is close to $\frac43$.
  \item Why is the result almost the same for every \omterm{def:g12:text-search:pattern}{pattern} length?
  \item With $26$ letters, what number replaces $\frac43$? Compare with the output of \code{compare_counts.py}.
\end{enumerate}

\needspace{4\baselineskip}\medskip\noindent\textbf{Part III --- Horspool's search.}
\begin{enumerate}[resume]
  \item The mean jump is the average of the \omterm{def:g12:text-search:shift-table}{shift table} over the four letters. Compute it for the \omterm{def:g12:text-search:pattern}{pattern}
    \code{ACGTACGT}.
  \item At each alignment Horspool makes a little more than one comparison. Explain why its count per letter is
    about $1.4$ divided by the mean jump, and check it on the output.
  \item Why does the mean jump grow from $m = 4$ to $m = 16$, and then stop growing?
  \item Why can no shift exceed $m$, and why, with four letters, are shifts above $10$ rare for long \omterm{def:g12:text-search:pattern}{patterns}?
  \item Which \omterm{def:g12:text-search:pattern}{pattern} of length $8$ has the largest mean jump? Give its \omterm{def:g12:text-search:shift-table}{shift table}.
\end{enumerate}

\needspace{4\baselineskip}\medskip\noindent\textbf{Part IV --- The ratio.}
\begin{enumerate}[resume]
  \item Read the ratio of the two counts for each $m$.
  \item Compare with the chart for $26$ letters: what would the ratio be at $m = 32$ there?
  \item The count for $m = 32$ is higher than for $m = 16$. Is the search worse for long \omterm{def:g12:text-search:pattern}{patterns}? What does the
    variation say about five \omterm{def:g12:text-search:pattern}{patterns} per length?
  \item The comparisons are counted, not timed. Which listing of the chapter shows that the two do not tell the same
    story, and why?
  \item What text and \omterm{def:g12:text-search:pattern}{pattern} would make Horspool's count per letter close to $m$?
  \item Why would a real program use \code{str.find} rather than either of these functions?
  \item State the result: comparisons per letter of the two searches on the chain, their ratio, and how it depends on
    the \omterm{def:g12:text-search:pattern}{pattern}'s length.
\end{enumerate}
\end{problem}
